% Home-screen / web-app icon, generated rather than drawn so it can be redone
% from scratch and stays consistent with the site's navbar colour.
%
% Regenerate (from this directory):
%   pdflatex icon.tex && sips -s format png icon.pdf --out icon-512.png
%   sips -Z 192 icon-512.png --out icon-192.png
%   sips -Z 180 icon-512.png --out apple-touch-icon.png
%
% sips rasterises PDF at 72dpi, so a 512bp page yields exactly 512px.
% Text is kept well inside the centre so Android's maskable crop cannot clip it.
\documentclass{article}
\usepackage[paperwidth=512bp,paperheight=512bp,margin=0bp]{geometry}
\usepackage{mathpazo}
\usepackage{xcolor}
\definecolor{tile}{HTML}{F6D85E}
\definecolor{tileink}{HTML}{222222}
\pagecolor{tile}
\pagestyle{empty}
\setlength{\parindent}{0pt}
\begin{document}
\vspace*{\fill}
\begin{center}
  \color{tileink}\bfseries
  {\fontsize{140bp}{150bp}\selectfont AI}\\[24bp]
  {\fontsize{140bp}{150bp}\selectfont Policy}
\end{center}
% Optical centring: the last line's text box reserves descender space that has
% no ink in these all-caps/digit lines, so the glyphs sit ~11bp low when the
% box is centred. This rigid space rebalances it (measured on the rendered PNG).
\vspace*{11bp}
\vspace*{\fill}
\end{document}
